% Methodik-Kapitel für die Bachelorarbeit
\chapter{Methodik}
\label{ch:Methodik}

\section{Forschungsdesign}
\label{ch:Methodik:sec:Forschungsdesign}

Ziel dieses Kapitels ist es, die methodische Vorgehensweise zu beschreiben, mit der modellbasierte Simulationen zur Untersuchung von Fahrzeugbewegungen und Weichensteuerung in dieser Arbeit durchgeführt werden. Die Arbeit verfolgt einen konstruktiv-experimentellen Forschungsansatz: Es wird eine softwarebasierte Simulationsumgebung entwickelt und genutzt, mit der gezielt experimentelle Fragestellungen geprüft werden. Die zentrale Forschungsfrage lautet, wie verschiedene Pfadfindungs- und Abstandsprüfstrategien in einem modularen Simulationsmodell das Verhalten von Fahrzeugen und das Gesamtsystem beeinflussen.

Die Wahl des Simulationsansatzes begründet sich durch folgende Anforderungen:
- Reproduzierbarkeit: Experimente müssen deterministisch konfigurierbar und wiederholbar sein.
- Modularität: Komponenten wie Streckenmodell, Fahrzeugmodell, Visualisierung und Strategien müssen unabhängig austauschbar sein.
- Messbarkeit: Die Umgebung muss Metriken (z. B. Laufzeit, Abstände, Kollisionen, Pfadlänge) automatisiert erfassen können.

Als Fallstudie dient das im vorliegenden Repository implementierte Simulationsprojekt. Wesentliche Implementierungsbereiche sind der Einstiegspunkt \texttt{main.py} sowie die Module \texttt{models/Track.py}, \texttt{vehicle.py}, \texttt{models/VehicleController.py} und die Pfadfindungsstrategien in \texttt{strategies/pathfinding/}. Die Experimente werden auf synthetisch generierten und gespeicherten Streckensätzen aus \texttt{saved\_tracks/} ausgeführt, um Variation und Vergleichbarkeit zu gewährleisten.

\subsection*{Versuchsplanung und Hypothesen}
- Hypothese 1: Unterschiedliche Pfadfindungsalgorithmen (z. B. Dijkstra vs. angepasste Varianten) führen messbar zu unterschiedlichen Laufzeiten und Pfadlängen.
- Hypothese 2: Erweiterte Abstandsprüfungen reduzieren Kollisionen auf Kosten der Fahrzeit.
- Unabhängige Variablen: gewählte Pfadfindungsstrategie, Abstandsprüfparameter, Track-Topologie.
- Abhängige Variablen: Anzahl Kollisionen, mittlere Fahrzeit, Pfadlänge, Rechenaufwand (Laufzeit).

Der Experimentaufbau folgt einem vollständig protokollierten Ablauf: (1) Auswahl bzw. Generierung eines Teststreckensets, (2) Initialisierung der Simulationsparameter, (3) Ausführung deterministischer Simulationen, (4) automatisierte Protokollierung der Resultate, (5) statistische Auswertung.

\section{Simulationsumgebung}
\label{ch:Methodik:sec:Simulationsumgebung}

\subsection*{Architekturüberblick}
Die Simulationsumgebung ist als modulare Softwarearchitektur implementiert. Die Hauptkomponenten sind:
- Eingangs-/Steuerdatei und Runner: \texttt{main.py} orchestriert Simulationen, lädt Streckendaten aus \texttt{saved\_tracks/} und initialisiert Szenarien.
- Modellschicht: In \texttt{models/} befinden sich Kernklassen wie \texttt{Track} (\texttt{models/Track.py}) und Segmentklassen (\texttt{models/segments/}), die topologische und geometrische Eigenschaften der Strecke repräsentieren.
- Fahrzeugschicht: \texttt{vehicle.py} modelliert Fahrzeugzustand und -dynamik; der Controller steuert Fahrbefehle und Interaktion mit Pfadfindung.
- Strategien: Unter \texttt{strategies/} sind Pfadfindungs- und Abstandsprüfalgorithmen getrennt abgelegt (z. B. \texttt{strategies/pathfinding/DijkstraRail.py}).
- Visualisierung/Monitoring: \texttt{view.py} und Module in \texttt{view/segments/} erzeugen Darstellungen und erlauben manuelle Inspektion; Rohdatenexport ermöglicht externe Auswertung.

\subsection*{Schichtenmodell}
Die Architektur lässt sich in vier Schichten gliedern:
- Persistenz- und Eingabeschicht: Laden/Speichern von Streckendefinitionen und Versuchskonfigurationen.
- Domänenschicht / Modell: Repräsentiert geometrische und logische Modelle (z. B. \texttt{models/Track.py}).
- Logik-/Strategieschicht: Pfadfindung und Abstandsprüfungen in \texttt{strategies/} sowie Steuerlogik im Controller.
- Präsentations- und Auswertungsschicht: Visualisierung durch \texttt{view.py} und Export/Logging von Messdaten.

\subsection*{Module und Verantwortlichkeiten}
- \texttt{models/Track.py}: Aufbau und Validierung der Strecke, Verbindung von Segmenten, Topologieabfragen.
- \texttt{models/segments/*}: Implementierung unterschiedlicher Segmenttypen (Kurve, Gerade, Weiche) inklusive Übergangsbedingungen.
- \texttt{vehicle.py}: Zustand, Bewegungsgesetz, Geschwindigkeits- und Bremsprofile.
- \texttt{VehicleController.py}: Schnittstelle zwischen Hochlevel-Strategie (Pfad) und Low-Level-Befehlen; behandelt Reservierungen, Blockaden und Weichensteuerung.
- \texttt{strategies/pathfinding/*}: Pfadplaner (z. B. Dijkstra-basierte Implementierung) und Schnittstelle für Erweiterungen.
- \texttt{strategies/distance\_checks/*}: Sicherheitsprüfungen und Abstandslogik, die Brems- und Ausweichverhalten beeinflussen.
- \texttt{view.py} und \texttt{view/segments/*}: Zeichnet Szenarien, unterstützt visuelles Debugging und Export.

\subsection*{Implementierungsdetails und technische Entscheidungen}
- Datenmodelle: Strecken und Segmente sind als serialisierbare JSON-Strukturen gehalten (Beispiele in \texttt{saved\_tracks/}), um Experimentkonfigurationen reproduzierbar zu machen.
- Determinismus: Zufallsquellen werden über Seeds steuerbar gehalten, damit Läufe vergleichbar sind.
- Mess-API: Die Simulationsschleife exponiert Hooks zum Erfassen von Metriken (Positions-Logs, Zeitstempel, Statusmeldungen). Ergebnisse werden in JSON/CSV geschrieben.
- Erweiterbarkeit: Die \texttt{strategies}-Schnittstelle erlaubt das Hinzufügen neuer Pfadfinder oder Abstandsprüfer ohne Änderung der Modell- oder Visualisierungsschicht.

\subsection*{Grafik und Visualisierung}
Zur Unterstützung der Interpretation werden folgende Visualisierungen eingesetzt:
- Zeitreihenplots (Geschwindigkeit, Abstand) pro Fahrzeug.
- Visualisierung der befahrenen Pfade überlagert mit Segmenttypen (Kurven, Weichen).
- Heatmaps zur Darstellung von Konfliktbereichen und Stopps.

Rohdaten werden in tabellarischer Form exportiert und können in Jupyter Notebooks oder mit Matplotlib weiter aufbereitet werden; die in \texttt{view.py} vorhandenen Funktionen erleichtern schnelle Inspektionen.

\section{Experimentelle Vorgehensweise}
\label{ch:Methodik:sec:Experimentelle Vorgehensweise}

\subsection*{Anforderungen an Testumgebungen}
Die Testumgebung muss folgende Anforderungen erfüllen:
- Wiederholbarkeit: Automatisierbare Skripte (z. B. Batch-Läufe via \texttt{main.py}) für deterministische Ausführungen.
- Messbarkeit: konsistente Zeitmessung, Fehler- und Ereignis-Logging sowie strukturierter Ergebnisexport.
- Skalierbarkeit: Möglichkeit, mehrere Szenarien sequenziell auszuführen und Ergebnisse zentral zu sammeln.
- Transparenz: alle Eingabeparameter und Seeds werden mit den Resultaten gespeichert, sodass jede Messung nachvollziehbar ist.

\subsection*{Konfigurationsmanagement}
Jeder Versuch wird durch eine Konfigurationsdatei beschrieben (JSON/YAML). Typische Parameter sind:
- Pfadfindungsstrategie und ihre Parameter.
- Abstandsprüfparameter (Mindestabstand, Reaktionszeit, Bremsprofile).
- Ausgewählte Streckendefinition (Dateien aus \texttt{saved\_tracks/}), Start-/Zielpaare.
- Simulationsdauer, Zeitschritt und Logging-Intervalle.

\subsection*{Versuchsvariationen}
Systematische Variationen ermöglichen das Testen der Hypothesen:
- Algorithmusvergleich: Identische Szenarien mit verschiedenen Pfadfindern (z. B. \texttt{DijkstraRail} vs. heuristische Varianten).
- Sicherheitsparameter: Variation von Mindestabständen und Bremsverzögerungen.
- Topologievariation: Nutzung unterschiedlicher Layouts aus \texttt{saved\_tracks/} (z. B. Ovale, Weichenkombinationen, enge Kurven).
- Lastszenarien: Unterschiedliche Fahrzeugdichten und gleichzeitige Startaufträge zur Erzeugung von Konflikten an Weichen.

\subsection*{Messgrößen und Auswertung}
Erfasste Metriken pro Experiment:
- Laufzeit der Simulation (reale Rechenzeit).
- Fahrzeit jedes Fahrzeugs (Zeit von Start bis Ziel).
- Pfadlänge und Anzahl der Weichenstellungen.
- Anzahl kritischer Ereignisse (Konflikte, Stopps, Kollisionen).
- Infrastrukturmetriken (Weichenstellungsanzahl, Blockadezeiten).

Statistische Auswertung umfasst mehrfache Durchläufe pro Konfiguration zur Berechnung von Mittelwerten, Standardabweichungen und ggf. Signifikanztests (z. B. t-Test oder nichtparametrische Alternativen). Visualisierungen (Boxplots, Zeitreihen, Heatmaps) unterstützen die Interpretation.

\subsection*{Ablauf eines einzelnen Experiments}
1. Auswahl der Konfiguration und der Track-Datei aus \texttt{saved\_tracks/}.
2. Initialisierung: Lade Strecke (\texttt{models/Track.py}), instanziere Fahrzeuge (\texttt{vehicle.py}) und Controller, setze Strategien.
3. Start: Starte Simulationslauf via \texttt{main.py} mit festgelegtem Seed.
4. Monitoring: Sammle Logs und Metriken während der Simulation.
5. Abschluss: Schreibe Resultate in strukturierte Dateien, generiere Basisplots.
6. Wiederholung: Wiederhole Schritte für alle gewünschten Variationen und Seeds.

\subsection*{Risiken, Limitationen und Maßnahmen}
- Modellvereinfachungen: Realistische Fahrdynamik ist nur approximiert; absolute Zeitwerte sind eher relativ zu interpretieren. Maßnahme: Fokus auf Relativvergleiche zwischen Strategien.
- Determinismus: Externe Systembelastung kann Rechenzeitmessungen beeinflussen. Maßnahme: Läufe auf ruhender Maschine durchführen und mehrere Wiederholungen verwenden.
- Komplexitätswachstum: Große Szenarien können Rechenzeit stark erhöhen. Maßnahme: Pilotläufe und adaptive Stichprobengrößen planen.

\section*{Zusammenfassung}
Die kombinierte Vorgehensweise — modularer Entwurf der Simulationsumgebung, strukturierte Experimentkonfigurationen und systematische Variationen — erlaubt robuste Vergleiche zwischen Pfadfindungs- und Abstandsstrategien. Die Implementierung im vorliegenden Projekt (u. a. \texttt{main.py}, \texttt{models/Track.py}, \texttt{vehicle.py}, \texttt{strategies/}) liefert die technische Basis, die in den folgenden Kapiteln zur Durchführung und Auswertung der Experimente genutzt wird.

% Ende Methodik
